
\textbf{RLHF Methods and Evaluation.} RLHF trains language models by learning reward functions from human preference pairs \citep{ouyang2022training}. The HH-RLHF dataset \citep{bai2022training} contains 161K preference pairs annotated for helpfulness and harmlessness. Standard benchmarks evaluate response quality and safety but do not measure effects on user agency. Shapira et al. (2026) identify sycophancy (over-agreement with users) as an emergent RLHF property, suggesting potential mechanisms by which preference optimization might affect agency-preserving language patterns.

\textbf{Linguistic Markers in Psychology.} Juanchich et al. (2017) demonstrate through controlled experiments that modal verbs and pronoun usage affect perceived autonomy attribution. Biber et al. (1999) document hedging markers as epistemic stance indicators in human corpora. These studies validate markers in human language but do not test transfer to AI-generated text.

\textbf{Bidirectional Alignment Framework.} Shen et al. (2024) systematically review 400+ papers and identify two alignment dimensions: $AI\rightarrow Human$ (integrating human specifications into systems) and $Human\rightarrow AI$ (preserving human agency). They note the absence of computational operationalization for the $Human\rightarrow AI$ dimension. This study tests one potential operationalization approach.
